\title{xLSTM: Extended Long Short-Term Memory}

\begin{abstract}
The 1990s saw the development of the constant error carousel and gating mechanisms, which became foundational to Long Short-Term Memory (LSTM) networks. LSTMs have demonstrated remarkable durability and enabled significant achievements in deep learning, notably forming the basis for early Large Language Models (LLMs). With the emergence of the Transformer architecture—featuring highly parallelizable self-attention—a new paradigm emerged, one that rapidly surpassed LSTMs at larger scales. This leads us to a key inquiry: What performance can be attained in language modeling if LSTMs are scaled up to billions of parameters, employing advanced methods inspired by contemporary LLMs while actively addressing their well-known shortcomings? To this end, we propose the use of exponential gating coupled with robust normalization and stabilization methodologies. We further refine the LSTM memory organization by designing: (i) sLSTM, which utilizes scalar memory, scalar updates, and an innovative approach to memory mixing, and (ii) mLSTM, which enables full parallelization via matrix-based memory and an update rule based on covariance. By embedding these upgraded LSTM variants within residual block frameworks, we construct xLSTM blocks that are further assembled through residual connections to produce xLSTM models. The introduction of exponential gating and enhanced memory mechanisms allow xLSTM to achieve competitive performance relative to advanced Transformers and State Space Models, demonstrating improvements in both accuracy and scaling behavior.\\
Code available at: \url{https://github.com/NX-AI/xlstm}
\end{abstract}

\section{Introduction}
\label{sec:intro}

The Long Short-Term Memory (LSTM) ideas~\citep{Hochreiter:91,Hochreiter:97nips,Hochreiter:97}, i.e., the constant error carousel and gating, were introduced to overcome the vanishing gradient problem of recurrent neural networks~\citep{Hochreiter:91,Hochreiter:00book}:

\begin{align}
\label{eq:lstm_idea}
  c_t \ = \ f_t \ c_{t-1} \ + \ 
          i_t \  z_t \ , \quad  
          h_t \ = \ o_t \ \psi({c_t}) \ . 
\end{align}

The cell state $c_{t-1}$ (depicted in green) undergoes an additive modification via the incorporation of cell inputs $z_t$, with this process regulated through sigmoid-based gates (blue). The forget gate $f_t$ and the input gate $i_t$ are responsible for governing the state update, whereas the output gate $o_t$ dictates the resulting memory cell output, specifically the hidden state $h_t$. Before producing the hidden state, the cell state is first normalized or compressed using the function $\psi$, after which the output gate modulates the final hidden state.

LSTMs have found wide-ranging success across multiple fields \citep{Hochreiter:01,Hochreiter:07,Schmidhuber:15}, holding dominance in text generation tasks up to the emergence of Transformers in 2017~\citep{Vaswani:17}. Their utility has been extensively validated in diverse sequence modeling tasks, such as natural language generation~\citep{Graves:13,Karpathy:15blog}, handwriting synthesis~\citep{Graves:13}, sequence-to-sequence machine translation~\citep{Sutskever:14nips}, program analysis~\citep{Zaremba:14arxiv}, image caption creation~\citep{Karpathy:15,Hossain:19}, code synthesis~\citep{Karpathy:15blog}, hydrological modeling (e.g., rainfall-runoff predictions)~\citep{Kratzert:18,Kratzert:19}, as well as flood warning systems~\citep{Nearing:24}. Within reinforcement learning, LSTM architectures have consistently delivered top performance, as demonstrated by implementations such as AlphaStar for StarCraft II~\citep{Vinyals:17short}, OpenAI Five for Dota 2~\citep{OpenAI:19}, and controllers developed for magnetic regulation in nuclear fusion~\citep{Degrave:22short}. LSTMs exhibit a remarkable ability to model abstractions, skillfully encoding semantic content within their cell state, evidenced by the identification of specific neurons specialized for number and syntactic structure~\citep{Lakretz:19}, language-related features~\citep{Bau:19}, or sentiment analysis~\citep{Radford:17arxiv}. Notwithstanding ongoing advancements in the field, LSTMs continue to underpin several crucial real-world systems~\citep{Degrave:22short,Nearing:24}, attesting to their enduring robustness and relevance.

Although LSTMs have achieved remarkable results, they exhibit three principal shortcomings: (i) Inability to alter prior storage decisions. We illustrate this drawback with the {\it Nearest Neighbor Search} scenario (refer also to Appendix~\ref{sec:appNearestNeighborSearch}): the model must process a sequence in search of the vector most akin to a given reference and, at the end, retrieve its associated value. The mean squared error for this task is depicted in the left panel of Figure~\ref{fig:lstmProblems}. When LSTM encounters a vector more similar than the one currently stored, it has difficulty replacing the original with the improved match; by contrast, our proposed xLSTM overcomes this via exponential gating. (ii) Restricted capacity for storage; LSTMs must condense all information into their scalar cell states. This bottleneck is evidenced in {\it Rare Token Prediction}, visualized in the right panel of Figure~\ref{fig:lstmProblems}, where the perplexity of predicting rare tokens from the Wikitext-103 corpus~\citep{Merity:17} is plotted across token frequency bins. LSTM displays poorer performance on less frequent tokens, attributable to its limited storage resources, whereas the xLSTM, with its matrix-based memory, significantly alleviates this issue. (iii) Inherent limitations to parallelization caused by hidden-to-hidden state dependencies that mandate sequential computation steps.

These challenges inherent to LSTM architectures have, in part, motivated the ascent of Transformers~\citep{Vaswani:17} in the domain of language modeling. It is therefore natural to ask: if we address these shortcomings and scale up LSTMs to the parameter regimes characteristic of modern Large Language Models, what performance levels can be realized in language modeling tasks?

\section{Extended Long Short-Term Memory}
\label{sec:xLSTM}

To address the constraints inherent in LSTM, the Extended Long Short-Term Memory (xLSTM) framework incorporates two principal changes to the original LSTM concept from Equation~\eqref{eq:lstm_idea}. These alterations—specifically, exponential gating and the design of innovative memory mechanisms—result in two expanded variants within the LSTM family: (i) the sLSTM (discussed in Section~\ref{sec:sLSTM}), which operates with scalar memory, applies scalar updates, and enables memory mixing; and (ii) the mLSTM (see Section~\ref{sec:mLSTM}), which utilizes a matrix-based memory together with a covariance-based (outer product) update, thereby allowing for full parallel processing. Both sLSTM and mLSTM capitalize on exponential gating to improve upon the standard LSTM. Importantly, mLSTM forgoes memory mixing—removing recurrent connections among hidden states—to facilitate parallelization. Both sLSTM and mLSTM are extendable to architectures with multiple memory cells; in sLSTM, memory mixing can occur across cells. Additionally, sLSTM is capable of supporting multiple heads, in which memory mixing happens among cells within each head but not between different heads. Incorporating heads into sLSTM, in conjunction with exponential gating, gives rise to an innovative approach to memory mixing. For mLSTM, the concepts of multiple heads and multiple memory cells are functionally indistinguishable.

Embedding these LSTM derivatives within residual block modules creates xLSTM blocks (further details in Section~\ref{sec:xLSTMblock}). Architectures constructed by stacking such xLSTM blocks residual-wise are termed xLSTM architectures, as outlined in Section~\ref{sec:xLSTMarchitecure}. The overall structure of an xLSTM architecture and its constituent components is illustrated in Figure~\ref{fig:xlstm_sketch}.

\subsection{Review of the Long Short-Term Memory}
\label{sec:orgLSTM}

The foundational concept behind LSTM~\citep{Hochreiter:91,Hochreiter:97nips,Hochreiter:97} centers on a scalar memory cell designed to serve as both processor and storage, mitigating the vanishing gradient problem~\citep{Hochreiter:91,Hochreiter:00book} via the constant error carousel mechanism associated with cell state updates. The memory cell operates with three distinct gates: input, output, and forget, where the forget gate was subsequently added by \citet{Gers:00}. The LSTM memory cell updates its state at each time step $t$ according to the following rules:

\begin{align}
c_t \ &= \  f_t \ c_{t-1} \ + \ i_t \ z_t &  & &\text{cell state} \\
h_t  \ &= \ o_t \ \tilde{h}_t \ , & \tilde{h}_t \ &= \ \psi \left( c_t\right)
  &\text{hidden state} \\
z_t \ &= \ \varphi \left( \tilde{z}_t \right) \ , 
  &\tilde{z}_t \ &=  \ \mathbf{w}^\top_{z} \ \mathbf{x}_t \ + \
  r_{z}  h_{t-1} \ + \  b_{z} \ \
  &\text{cell input} \\
i_t \ &= \ \sigma \left( \tilde{i}_t  \right) \ , 
  &\tilde{i}_t \ &= \ \mathbf{w}^\top_{i} \ \mathbf{x}_t \ + \
  r_{i}  \ h_{t-1} \ + \  b_{i} \ \
  &\text{input gate} \\
f_t \ &= \ \sigma \left(  \tilde{f}_t \right) \ , 
  &\tilde{f}_t \ &= \ \mathbf{w}^\top_{f} \ \mathbf{x}_t  \ + \
  r_{f}  \ h_{t-1} \ + \  b_{f} \ \
  &\text{forget gate} \\
o_t \ &= \ \sigma \left( \tilde{o}_t \right) \ , 
  &\tilde{o}_t  \ &= \ \mathbf{w}^\top_{o} \ \mathbf{x}_t \ + \
  r_{o}  \ h_{t-1} \ + \  b_{o} \ \
  &\text{output gate} 
\end{align}

The vectors $\mathbf{w}_{z}$, $\mathbf{w}_{i}$,
$\mathbf{w}_{f}$, and $\mathbf{w}_{o}$ denote the input weights linking the input $\mathbf{x}_t$ to the cell input, input gate, forget gate, and output gate, respectively. The parameters $r_{z}$, $r_{i}$, $r_{f}$, and $r_{o}$ represent the recurrent weights from the previous hidden state $h_{t-1}$ to the cell input, input gate, forget gate, and output gate. The bias variables $b_{z}$, $b_{i}$, $b_{f}$, and $b_{o}$ are included for each corresponding part. The functions $\varphi$ and $\psi$ act as the activation functions for the cell input and hidden state, commonly taken as $\tanh$. The $\psi$ activation functions to constrain the cell state within bounds; otherwise, the cell state could diverge. The gates employ the sigmoid activation function, defined by $\sigma \left( x \right)= 1/(1+\exp(-x))$. Subsequent modifications extended to vectors of scalar memory cells, $\mathbf{c}_t \in \mathbb{R}^{d}$, enabling recurrent weight matrices $\mathbf{R} \in \mathbb{R}^{d \times d}$ to interact and mix outputs among memory cells~\citep{Greff:15}, as further discussed in Appendix~\ref{sec:apporgLSTM}. Experimental studies indicate that every element of the memory cell is essential for performance~\citep{Greff:15}.

\subsection{sLSTM}
\label{sec:sLSTM}

In order to equip LSTMs with the capacity to modify their storage choices, we incorporate exponential gates (highlighted in red), along with mechanisms for normalization and stabilization. Specifically, both the input and forget gates may utilize exponential activation functions. To achieve normalization, we propose a normalizer state that accumulates the sum of the input gate values multiplied by all subsequent forget gate activations.

The scalar sLSTM forward pass is:

\begin{align}
\label{eq:slstmforward}
c_t \ &= \  f_t \ c_{t-1} \ + \ i_t \ z_t & & 
 &\text{cell state} \\
n_t \ &= \  f_t \ n_{t-1} \ + \ i_t  & & 
 &\text{normalizer state} \\
h_t \ &= \ o_t \ \tilde{h}_t \ ,
  &\tilde{h}_t \ &= \ c_t / n_t
&\text{hidden state} \\
z_t \ &= \ \varphi \left( \tilde{z}_t \right) \ , 
  &\tilde{z}_t \ &=  \ \mathbf{w}^\top_{z} \ \mathbf{x}_t \ + \
  r_{z}  h_{t-1} \ + \  b_{z}
  &\text{cell input} \\
i_t \ &= \ \!\exp \left( \tilde{i}_t  \right) \ , 
  &\tilde{i}_t \ &= \ \mathbf{w}^\top_{i} \ \mathbf{x}_t \ + \
  r_{i}  \ h_{t-1} \ + \  b_{i}
  &\text{input gate} \\
f_t \ &= \ \sigma \left(  \tilde{f}_t \right) \ \text{OR} \
\!\exp \left(  \tilde{f}_t \right) \ ,
  &\tilde{f}_t \ &= \ \mathbf{w}^\top_{f} \ \mathbf{x}_t  \ + \
  r_{f}  \ h_{t-1} \ + \  b_{f}
  &\text{forget gate} \\
o_t \ &= \ \sigma \left( \tilde{o}_t \right) \ , 
  &\tilde{o}_t  \ &= \ \mathbf{w}^\top_{o} \ \mathbf{x}_t \ + \
  r_{o}  \ h_{t-1} \ + \  b_{o}
  &\text{output gate} 
\end{align}

We transfer the original LSTM gating techniques, i.e., input- and/or hidden-dependent gating plus bias term, to the new architectures. Exponential activation functions can lead to large values that cause overflows. Therefore, we stabilize gates with an additional state \mbox{$m_t$~\citep{Milakov:18arxiv}}:

\begin{align}
\label{eq:slstmstabil}
m_t \ &= \ \max \left( \log ( f_t ) + m_{t-1} , \log ( i_t ) \right) &\text{stabilizer state} \\
i'_t \ &= \ \!\exp \left( \log \left ( i_t \right) - m_t \right) \ = \ \!\exp \left( \tilde{i}_t - m_t \right) \ \, 
  &\text{stabil. input gate} \\
f'_t \ &= \ \!\exp \left( \log \left( f_t \right) + m_{t-1} - m_t \right) \ \, 
  &\text{stabil. forget gate}
\end{align}

As demonstrated in Appendix~\ref{sec:appsLSTM}, substituting $f_t$ with $f'_t$ and $i_t$ with $i'_t$ during the forward computation leaves both the network’s output and the gradients of the loss with respect to the parameters unchanged.

\paragraph{New Memory Mixing.} Similar to conventional LSTM architectures, sLSTM can support several memory cells (refer to Appendix~\ref{sec:appsLSTM}). The presence of multiple memory cells facilitates memory mixing through the recurrent mappings $\mathbf{R}_{\mathbf{z}}$, $\mathbf{R}_{\mathbf{i}}$, $\mathbf{R}_{\mathbf{f}}$, and $\mathbf{R}_{\mathbf{o}}$ that connect the hidden state vector $\mathbf{h}$ to the memory cell input $\mathbf{z}$ as well as the gates $\mathbf{i}$, $\mathbf{f}$, and $\mathbf{o}$. A novel feature in memory mixing emerges from exponential gating. With the redesigned sLSTM, each head allows internal memory mixing while interactions between heads are absent. This addition of multiple heads in sLSTM, paired with exponential gating, introduces a fundamentally new mechanism for memory mixing.

\subsection{mLSTM}
\label{sec:mLSTM}

To augment the memory capacity of LSTMs, we replace the conventional scalar cell state $c \in \mathbb{R}$ with a matrix cell state $\mathbf{C} \in \mathbb{R}^{d \times d}$. Consequently, memory access now involves matrix multiplication operations. At each time step $t$, we aim to store a vector pair, consisting of a key $\mathbf{k}_t \in \mathbb{R}^d$ and a value $\mathbf{v}_t \in \mathbb{R}^d$—following terminology from the Transformer framework. Later, at a future time $t + \tau$, the value $\mathbf{v}_t$ can be recovered using a query vector $\mathbf{q}_{t+\tau} \in \mathbb{R}^d$. This setup aligns with the framework of Bidirectional Associative Memories (BAMs) as described in \citep{Kohonen:72,Anderson:72,Nakano:72,Anderson:77}.

The covariance update rule~\citep{Sejnowski:77,Dayan:91} for storing a key-value pair is

\begin{align}
\mathbf{C}_t \ &= \ \mathbf{C}_{t-1} \ + \ \mathbf{v}_{t} \ \mathbf{k}_t^\top \ .
\end{align}

We consider the inputs undergo layer normalization prior to their projection into key and value representations, ensuring the mean is zero. The covariance-based update mechanism is provably optimal~\citep{Dayan:91} for achieving maximal separability among binary vectors being retrieved, which corresponds to optimizing the signal-to-noise ratio. By restricting retrieval to only pairwise terms and accepting the quadratic computational burden—as in attention mechanisms~\citep{Krotov:16,Krotov:17,Ramsauer:21}—even greater separability can be obtained. The covariance update formulation is also synonymous with the Fast Weight Programmers paradigm~\citep{Schmidhuber:92ncfastweights,Schlag:21}; in later extensions, these have been supplemented with a fixed decay parameter for $\mathbf{C}_{t-1}$ and a fixed learning rate applied to $\mathbf{v}_{t} \mathbf{k}_t ^\top$~\citep{Ba:16}. Inspired by these advances, we incorporate the covariance update mechanism within the LSTM architecture, associating the forget gate with the decay parameter, the input gate with the learning rate, and assigning scaling of the retrieved vector to the output gate.

For this matrix memory, the normalizer state is the weighted sum of key vectors, where each key vector is weighted by the input gate and all future forget gates. Again, the normalizer state keeps record of the strength of the gates. Since the dot product between query and normalizer state can be close to zero, we use the absolute value of this dot product and lower bound it by a threshold (typically 1.0) as done previously~\citep{Sun:23arxiv}. The mLSTM forward pass is:

\begin{align}
\label{eq:mlstm_recurrent_begin}
\mathbf{C}_t \ &= \  f_t \ \mathbf{C}_{t-1} \ + \ 
  i_t \ \mathbf{v}_t \ \mathbf{k}_t^\top & &  &\text{cell state} \\
\mathbf{n}_t \ &= \  f_t \ \mathbf{n}_{t-1} \ + \ 
  i_t \ \mathbf{k}_t & &  &\text{normalizer state} \\
\mathbf{h}_t  \ &= \ \mathbf{o}_t \ \odot \ \tilde{\mathbf{h}}_t
\ , \qquad \qquad 
\tilde{\mathbf{h}}_t \ = \   \mathbf{C}_t \mathbf{q}_t \ / \ 
  \max \left\{ |\mathbf{n}_t^\top \mathbf{q}_t|, 1 \right\} 
   &&&\text{hidden state} \\
\mathbf{q}_t \ &= \ \mathbf{W}_q \ \mathbf{x}_t \ + \ \mathbf{b}_q  & &  &\text{query input} \\
\mathbf{k}_t \ &= \ \frac{1}{\sqrt{d}} \mathbf{W}_k \ \mathbf{x}_t \ + \ \mathbf{b}_k  & &  &\text{key input} \\
\mathbf{v}_t \ &= \ \mathbf{W}_v \ \mathbf{x}_t \ + \ \mathbf{b}_v  & &  &\text{value input} \\
i_t \ &= \ \!\exp \!  \! \left( \tilde{i}_t  \right) \ , 
 \qquad \qquad \ \ \ \ \,
  \tilde{i}_t \ = \ \mathbf{w}^\top_{i} \ \mathbf{x}_t \ + \  b_{i} \ \ \ 
  &&&\text{input gate} \\
f_t \ &= \ \sigma \!  \left(  \tilde{f}_t \right) \ \text{OR} \ \!\exp \!  \! \left(  \tilde{f}_t \right) , \ \ \: \!
\tilde{f}_t \ = \ \mathbf{w}^\top_{f} \ \mathbf{x}_t  \ + \
  b_{f}
  &&& \text{forget gate} \\
\label{eq:mlstm_recurrent_end}
\mathbf{o}_t \ &= \ \sigma \left( \tilde{\mathbf{o}}_t \right) \ , \qquad \qquad \quad \ \ \ \,
  \tilde{\mathbf{o}}_t  \ = \ \mathbf{W}_{\mathbf{o}} \ \mathbf{x}_t \ + \
  \mathbf{b}_{\mathbf{o}}
  &&&\text{output gate} 
\end{align}

mLSTM is capable of supporting several memory cells, similar to the conventional LSTM architecture. In the context of mLSTM, having multiple heads functions identically to having multiple cells, given the absence of memory mixing mechanisms. To maintain stability in the exponential gating functions within mLSTM, we apply the same stabilization approaches as those introduced for sLSTM (refer to Equation~\ref{eq:slstmstabil}). Due to the lack of memory interaction across cells in mLSTM, the recurrence relation admits a parallelizable formulation. Additional information can be found in Appendix~\ref{sec:appmLSTM}.

\subsection{xLSTM Architecture}
\label{sec:xLSTMarch}

\paragraph{xLSTM Blocks.}
\label{sec:xLSTMblock}
The xLSTM block is designed to produce a non-linear summary of previous information within a high-dimensional representation, thereby enhancing the differentiation between varying contexts or sequences. Proper distinction between these histories is essential for accurate prediction of subsequent sequence elements, such as the next token. Drawing on Cover’s Theorem~\citep{Cover:65}, which asserts that patterns that undergo non-linear embedding into higher dimensional spaces have a greater likelihood of being separated linearly compared to patterns in lower dimensional spaces, we explore two types of residual block structures: (i) The first structure, resembling the approach used in Transformers, implements a residual block where an up-projection is performed after a non-linear summary of past states in the original space; this is followed by a linear projection into a higher-dimensional space, application of a non-linear activation, and a subsequent linear projection back to the initial space, as shown in the left segment of Figure~\ref{fig:Backbones} and the third column of Figure~\ref{fig:xlstm_sketch}. Additional details are provided in Figure~\ref{fig:appsLSTM_detailed} of the appendix. (ii) The second architecture, akin to that of State Space Models, executes a linear projection into a high-dimensional space prior to further processing,

The process involves a non-linear compression of past information within a high-dimensional latent space, followed by a linear transformation that projects data back into its original space. When the xLSTM block incorporates an sLSTM, the post up-projection module is applied; conversely, for xLSTM blocks equipped with an mLSTM, the pre up-projection module is employed to leverage the expanded memory afforded by the high-dimensional context. For further illustration, please consult the left side of Figure~\ref{fig:Backbones}, the third column in Figure~\ref{fig:xlstm_sketch}, or Figure~\ref{fig:appsLSTM_detailed} within the appendix.

\paragraph{xLSTM Architecture. }
\label{sec:xLSTMarchitecure}
The xLSTM framework is assembled through the residual stacking of foundational modules, as established in prior work~\citep{Srivastava:15,He:16}. Our approach utilizes the prevalent pre-LayerNorm~\citep{Ba:16b} residual backbone architecture, mirroring those found in state-of-the-art Large Language Models. Refer to the last column in Figure~\ref{fig:xlstm_sketch} for a detailed depiction.

\subsection{Memory and Speed Considerations}

Unlike Transformers, xLSTM architectures feature linear computational complexity and maintain fixed memory usage regardless of input sequence length. The compressive nature of the xLSTM memory makes it highly appropriate for deployment in industrial settings and edge computing scenarios.

In mLSTM, memory operations are parameter-free but incur significant computational demands due to its reliance on $d \times d$ matrix representations and updates. This architecture prioritizes memory capacity at the expense of increased compute. Despite this, the workload is amenable to efficient parallelization on GPUs, causing only a negligible increase in wall clock runtime.

Although mLSTM can be efficiently parallelized in a manner similar to FlashAttention~\citep{Dao:22,Dao:23} and GLA~\citep{Yang:23arxiv}, sLSTM does not support such parallelization because of its memory mixing through hidden-to-hidden connections. Nonetheless, we have implemented a high-performance CUDA version, optimizing GPU memory usage down to the register level. As a result, sLSTM typically achieves runtimes less than twice as slow as mLSTM.

\section{Related Work}
\label{sec:related}

\paragraph{Linear Attention.} Numerous approaches have been proposed to mitigate the quadratic computational costs associated with Transformer models, aiming to make attention operations scale linearly with respect to context length. Synthesizer learns artificial attention weights independently of token-to-token interactions~\citep{Tay:20arxiv}. Linformer leverages a low-rank approximation for the self-attention matrix and further simplifies computation through linearization~\citep{Wang:20arxiv}. Linear Transformer implements a linearized form of the attention mechanism \citep{Katharopoulos:20}. Performer employs a positive orthogonal random features strategy to achieve a linear approximation of the attention softmax~\citep{Choromanski:21}. In addition, attention components have been supplanted by efficient, long-range convolutional operations in models such as Structured Global Convolution (SGConv)~\citep{Li:22} and Hyena Hierarchy~\citep{Poli:23}.

\paragraph{State Space Models.} State Space Models (SSMs) have gained considerable attention in recent times due to their linear scalability with context length and their competitive results relative to Transformer architectures. Among the initial contributions was the Structured State Space sequence model (S4)~\citep{Gu:21}, which was succeeded by models such as the Diagonal State Space (DSS)~\citep{Gupta:22}, Gated State Space (GSS)~\citep{Mehta:22}, S5~\citep{Smith:22}, Bidirectional Gated SSM (BiGS)~\citep{Wang:22}, H3~\citep{Fu:23}, and Mamba~\citep{Gu:24arxiv}.

\paragraph{Recurrent Neural Networks.} In recent work, Recurrent Neural Networks (RNNs) have been proposed as alternatives to Transformers and attention mechanisms, motivated by their linear scalability with respect to input sequence length. Notably, language modeling efforts with Deep Linear Recurrent Units (LRUs) have demonstrated strong results~\citep{Orvieto:23, De:24arxiv}, alongside methods such as Hierarchically Gated Linear RNN (HGRN)~\citep{Qin:23} and its subsequent version, HGRN2~\citep{Qin:24arxiv}. The RWKV model~\citep{Peng:23arxivshort,Peng:24arxivshort}, a prominent RNN-based approach for large-scale language modeling, achieves performance on par with Transformer-based architectures.

\paragraph{Gating.}
Gating, a fundamental concept introduced in LSTM architectures, has seen renewed attention and various reinterpretations in contemporary research. This mechanism has been integrated into numerous models, including HGRN~\citep{Qin:23}, HGRN2~\citep{Qin:24arxiv}, Gated Linear Attention (GLA)~\citep{Yang:23arxiv}, Gated State Space (GSS) models~\citep{Mehta:22}, Bidirectional Gated SSM (BiGS)~\citep{Wang:22}, Moving Average Equipped Gated Attention (MEGA)~\citep{Ma:22}, RWKV~\citep{Peng:23arxivshort}, and Mamba~\citep{Gu:24arxiv}.

\paragraph{Covariance Update Rule.}  
To boost memory capacity, we augmented the mLSTM cell with a matrix memory utilizing a covariance-based update rule. Related methods employing similar update techniques include Fast Weight Programmers \citep{Schmidhuber:92ncfastweights,Schlag:21}, RWKV-5 and RWKV-6~\citep{Peng:24arxivshort}, Retention~\citep{Sun:23arxiv}, Linear Transformer~\citep{Katharopoulos:20}, as well as HGRN2~\citep{Qin:24arxiv}.

\paragraph{Most Related.} The xLSTM architecture shares significant conceptual similarities with models such as Retention~\citep{Sun:23arxiv}, RWKV~\citep{Peng:23arxivshort,Peng:24arxivshort}, and HGRN2~\citep{Qin:24arxiv}, particularly regarding their use of matrix memory and gating mechanisms. Nonetheless, unlike the innovative sLSTM, these architectures do not incorporate memory mixing capabilities. The ability to mix memories is crucial for handling state tracking tasks, making LSTM variants more powerful than both State Space Models (SSMs) and Transformers~\citep{Merrill:24,Deletang:23} in such contexts. State tracking is essential, for instance, when executing code or monitoring entities throughout extensive narratives.

\paragraph{Residually Stacking Architectures.} The xLSTM architectures, akin to most modern large-scale deep learning models, are designed using the principle of residually stacking modular components~\citep{Srivastava:15,He:16}.

This approach has facilitated the development of deep convolutional neural networks~\citep{He:16} and forms the foundational structure of Transformers~\citep{Vaswani:17}. Transformers, in turn, have emerged as the driving mechanism powering Large Language Models (LLMs) such as GPT-3~\citep{Brown:20short}, ChatGPT~\citep{Schulman:22short}, GPT-4~\citep{Achiam:23gpt4short}, Megatron-LM~\citep{Shoeybi:19}, Gopher~\citep{Rae:21short}, ERNIE 3.0 Titan~\citep{Wang:21arxivshort}, GLaM~\citep{Du:21glamshort}, Chinese M6~\citep{Lin:21}, AlexaTM 20B~\citep{Soltan:22}, OPT~\citep{Zhang:22opt}, Chinchilla~\citep{Hoffmann:22short}, BLOOM~\citep{Scao:22arxivshort}, GLM-130B~\citep{Zeng:22arxivshort}, LaMDA~\citep{Thoppilan:22short}, PaLM~\citep{Chowdhery:22short}, Llama~\citep{Touvron:23arxiv}, as well as Gemini~\citep{GeminiTeam:23arxiv,Reid:24arxiv}.

\section{Experiments}
\label{sec:Exp}

In this section, we carry out empirical studies on xLSTM and benchmark it against existing approaches, with an emphasis on language modeling tasks. Section~\ref{sec:ExpSyntethic} explores the distinct features of xLSTM by applying it to synthetic task scenarios. Section~\ref{sec:ExpComparison} presents a comparison of validation perplexities across several state-of-the-art language modeling techniques, each trained with 15B tokens from the SlimPajama dataset~\citep{Soboleva:23}. Ablation studies are performed for xLSTM on this same dataset, and we analyze the scaling trends in line with methodologies from \citet{Kaplan:20} and \citet{Brown:20short}.

Section~\ref{sec:ExpLanguage} offers a comprehensive language modeling analysis: xLSTM and top-performing methods (as identified in Section~\ref{sec:ExpComparison}) are trained on a substantially larger corpus of 300B tokens from SlimPajama~\citep{Soboleva:23}. The evaluation includes: (i) measuring the methods’ ability to extrapolate to extended context lengths; (ii) determining their performance via validation perplexity and effectiveness on downstream tasks~\citep{Sutawika:24}; (iii) examining performance over 571 text domains sampled from the PALOMA language benchmark~\citep{Magnusson:23arxivshort}; and (iv) revisiting scaling properties, now with 20 times more training data than prior experiments.

\label{sec:xlstm_blocks_notation}
Throughout all experimental setups, we denote xLSTM[$a$:$b$] as representing the proportion $a/b$ of mLSTM-based to sLSTM-based xLSTM blocks. For instance, xLSTM[7:1] indicates that within eight blocks, seven are constructed with mLSTM and one utilizes sLSTM. When employing the standard block total of 48, this configuration yields 42 blocks using mLSTM and 6 using sLSTM. Additionally, in every experiment, we include both pre up-projection blocks for mLSTM and post up-projection blocks for sLSTM.

\subsection{Synthetic Tasks and Long Range Arena}
\label{sec:ExpSyntethic}
Initially, we evaluate how xLSTM's innovative exponential gating with memory mixing performs on formal language benchmarks~\citep{Deletang:23}. Subsequently, we investigate the advantages of xLSTM's matrix memory in the context of the Multi-Query Associative Recall task~\citep{Arora:23arxiv}. Lastly, we analyze xLSTM's ability to handle lengthy sequences by measuring its results on the Long Range Arena~\citep{Tay:21}.

\paragraph{Evaluation of xLSTM’s Exponential Gating and Memory Mixing Mechanism.} We investigate xLSTM’s innovative exponential gating combined with memory mixing, a feature purported to equip it for handling state tracking tasks~\citep{Merrill:24, Merrill:23}. To facilitate length extrapolation, we adopt and modify the formal language evaluation suite from \citet{Deletang:23} to allow training on sequences of varying lengths. Further details, including comprehensive task specifications and extended findings, are documented in Appendix~\ref{sec:appExpSynthetic-formalLang}. 

Our experimental setup contrasts xLSTM’s performance against alternative architectures such as Transformers, State Space Models, and traditional Recurrent Neural Networks. The models' accuracy is assessed exclusively on task-relevant tokens, with scores normalized so that 0 corresponds to random chance and 1 denotes flawless prediction. Specifically, we assess the following 2-block configurations: xLSTM[0:1] (representing sLSTM only), xLSTM[1:0] (representing mLSTM only), xLSTM[1:1], Llama, Mamba, RWKV, Retention, Hyena, LSTM, and Transformer block-integrated LSTM (LSTM (Block)). The outcomes are illustrated in Figure~\ref{fig:formal-main}. 

We observe that models lacking memory mixing, such as standard Transformers or State Space Models—both of which do not perform state tracking—are unable to solve basic regular grammar tasks like the parity challenge. These findings corroborate prior work indicating that Transformers and State Space Models are inherently less expressive than Recurrent Neural Networks for such stateful tasks~\citep{Merrill:24, Merrill:23, Deletang:23}.

\paragraph{Evaluation of xLSTM Memory on Associative Recall Benchmarks.} This study investigates the capabilities of the xLSTM’s novel matrix memory with respect to memory capacity, utilizing the Multi-Query Associative Recall task~\citep{Arora:23arxiv}. Each input sequence consists of key-value pairs, selected at random from an extensive vocabulary, which the model must retain for subsequent retrieval. To make this task more challenging than its original formulation, we scale the number of key-value pairs up to 256 and extend the possible sequence context to 2048. This adjustment provides a more comprehensive evaluation of various models' memory capacities. The comparison involves 2-block configurations of Llama, Mamba, RWKV-5, RWKV-6, xLSTM[1:1], and xLSTM[1:0]. Model performance is measured as the proportion of correctly recalled pairs. Due to the exponential relationship between memory and coding dimension in Transformers (such as Llama)~\citep{Ramsauer:21}, these architectures set a performance benchmark on this task. Figure~\ref{fig:mqar-main} depicts the main findings. Notably, xLSTM[1:1] achieves the highest accuracy among non-Transformer contenders, a pattern observable even for relatively compact model sizes. Contrary to expectations, incorporating the sLSTM block does not curtail memory capacity but in fact appears to enhance it, particularly evident at the hardest level of the benchmark featuring 256 pairs. Additional experimental details and extrapolation studies are provided in Appendix~\ref{sec:appExpSynthetic-mqar}, demonstrating that xLSTM’s memory improvements facilitate generalization to contexts beyond those encountered during training.

\paragraph{Test of xLSTM's Long Context Capabilities on Long Range Arena.} In order to evaluate xLSTM's ability to manage lengthy sequences and extended contexts, we conduct a comparative analysis of various approaches using the Long Range Arena~\citep{Tay:21}. Across all evaluated tasks, xLSTM delivers robust and reliable results, indicating that its architecture is highly effective and well-suited for tackling the diverse challenges associated with long context processing.

For more details, see Appendix~\ref{sec:appExpSynthetic-lra}.

\subsection{Method Comparison and Ablation Study}
\label{sec:ExpComparison}

In order to investigate the central research question—namely, the performance of our novel LSTM variants at scale for language modeling—we conduct experiments by training xLSTMs, Transformers, State Space Models, among other techniques, using 15B tokens sourced from SlimPajama under a consistent auto-regressive paradigm. The resulting models are evaluated on the validation dataset, and we carry out ablation analyses specifically for the xLSTM architectures.

\paragraph{Comparing xLSTM to Other Methods.} In this comparison, we train a suite of models using a 15B token subset from SlimPajama~\citep{Soboleva:23}. Each model’s effectiveness is assessed in terms of perplexity on a reserved validation set. The evaluated set of approaches includes: xLSTM (our contribution), GPT-3 (Transformer)~\citep{Brown:20short}, Llama (Transformer)~\citep{Touvron:23arxiv}, H3 (SSM)~\citep{Fu:23}, Mamba (SSM)~\citep{Gu:24arxiv}, RWKV-4, RWKV-5, RWKV-6 (RNNs)~\citep{Peng:23arxivshort,Peng:24arxivshort}, GLA (linear Transformer)~\citep{Yang:23arxiv}, HGRN and HGRN2 (RNNs)~\citep{Qin:23,Qin:24arxiv}, RetNet (linear Transformer)~\citep{Sun:23arxiv}, Hyena (linear Transformer)~\citep{Poli:23}, as well as variants xLSTM[1:0] and xLSTM[7:1] (further information in Section~\ref{sec:xlstm_blocks_notation}). Training was performed using mixed-precision arithmetic. Notably, for RWKV-5, RWKV-6, GLA, and HGRN2, mixed-precision computation did not leverage PyTorch’s automated mixed precision tools (refer to Appendix Section~\ref{sec:appGenTrainProc} for more details).

We partition these models into three categories: (a) Transformers, (b) State Space Models (SSMs), and (c) Recurrent Neural Networks (RNNs), with linear Transformers included alongside RNNs due to their use of linear mechanisms in place of standard Transformer attention. All models are designed to match the architecture size of GPT-3 with 350M parameters—specifically, embedding dimension of 1024 and 24 residual blocks. It is noteworthy that only GPT-3 implements weight sharing between token and output embeddings, resulting in a slightly reduced parameter count.

The data in Table~\ref{tab:spaj15b_model_results} demonstrate that xLSTM yields lower perplexity on the validation set than any previously existing alternative. Additional experiment specifics are available in Appendix~\ref{sec:appExpComparison}. Furthermore, Figure~\ref{fig:temp_sclaw15B} presents the scaling trends from these runs, suggesting the strong relative performance of xLSTM is likely to persist for models of greater scale.

\paragraph{Ablation Studies.}
Table~\ref{tab:spaj15b_model_results} and Figure~\ref{fig:temp_sclaw15B} illustrate that xLSTM attains outstanding performance on language modeling tasks after training with 15B tokens from SlimPajama. To systematically analyze the transition from LSTM to xLSTM, we incrementally adapt a standard LSTM model towards the xLSTM design. The process begins by incorporating LSTM layers within pre-LayerNorm residual frameworks. Next, the architecture is modified to include a post up-projection module. The final step involves introducing exponential gating mechanisms and matrix memory.

The results are shown in Table~\ref{tab:ablstudies} (top). 

The results of the ablation studies indicate that the observed improvements in performance are a consequence of the combined effects of exponential gating and matrix memory. Furthermore, considering the crucial role of gating within RNNs and State Space Models, we systematically examine various gating strategies. Table~\ref{tab:ablstudies} (bottom) demonstrates that allowing each gate to be learnable and modulated by the input contributes incremental gains. Further analyses regarding the isolated contributions of backbone components are provided in Appendix~\ref{sec:appAblDetails}.

\subsection{xLSTM as Large Language Model}
\label{sec:ExpLanguage}

Our investigation culminates in extensive language modeling experiments aimed at evaluating xLSTM as a large language model (LLM). To this end, we substantially scale up the training data and utilize 300B tokens drawn from SlimPajama—the same training volume adopted by prior works such as Mamba~\citep{Gu:24arxiv} and Griffin~\citep{De:24arxiv}.

In our comparisons, xLSTM is benchmarked against RWKV-4, Llama, and Mamba—each serving as the representative of a distinct methodological category described in Section~\ref{sec:ExpComparison}. RWKV-4 is chosen as the RNN baseline because appropriate precision settings for newer versions, namely RWKV-5, RWKV-6, and HGRN2, were established only after 300B-token training runs had commenced (for further details, consult Appendix~\ref{sec:appGenTrainProc}).

We investigate multiple model scales (125M, 350M, 760M, 1.3B parameters), assess the models’ capacity to extrapolate to longer contexts, and evaluate their efficacy on validation data. Furthermore, we analyze downstream task performance, measure their language modeling abilities across the 571 textual domains included in the PALOMA benchmark, and finally explore how their behavior aligns with established scaling laws.

\paragraph{Sequence Length Extrapolation.}
\label{sec:spaj300B_ctx_extrapolation}
We begin by assessing the ability of 1.3B parameter models—including xLSTM, RWKV-4, Llama, and Mamba—to extrapolate to longer sequence lengths. Each model is trained using a context window of 2048 tokens and subsequently evaluated on sequences as long as 16384 tokens. The outcomes are summarized in Figure~\ref{fig:spaj300B_seq_extrapolation}. Remarkably, xLSTM models are able to sustain low perplexity scores at extended context lengths, outperforming alternative architectures in this regard.

\paragraph{Validation Perplexity and Downstream Tasks.}
\label{sec:spaj_downstream_eval}
Next, across a range of model sizes, we compare the efficacy of xLSTM, RWKV-4, Llama, and Mamba architectures using two benchmarks: next token prediction accuracy on the SlimPajama validation set, and a suite of downstream tasks focused on common sense reasoning.

In Table~\ref{tab:lmeval}, the third column presents the perplexity values on the validation set for various approaches. Across every model size, xLSTM[1:0] and xLSTM[7:1] exhibit the lowest validation set perplexity, identifying them as the leading models in this metric. The remaining columns of Table~\ref{tab:lmeval} showcase results on downstream tasks. xLSTM consistently outperforms alternative methods in nearly all tasks and for every model size; the only exception is found in the ARC task, where Mamba sometimes achieves superior results. Additional discussion is included in Appendix~\ref{sec:appExpLanguage}.

\paragraph{Performance on PALOMA Language Tasks.} Thirdly, across all tested model scales, we evaluate xLSTM, RWKV-4, Llama, and Mamba models on PALOMA language tasks~\citep{Magnusson:23arxivshort}, focusing on their next token prediction abilities. We report perplexity scores for predicting subsequent tokens across 571 distinct text domains, spanning sources such as nytimes.com to subreddit communities like r/depression.

Table~\ref{tab:ppleval} provides perplexity results, organized into language modeling tasks (first seven columns) and detailed domain-specific benchmarks (last five columns). In these language evaluations, xLSTM[1:0] consistently yields better outcomes compared to xLSTM[7:1].

Examining model comparisons, xLSTM[1:0] achieves lower perplexity scores than Mamba in 568 of 571 domains (99.5\%), outperforms Llama in 486 domains (85.1\%), and surpasses RWKV-4 in 570 domains (99.8\%). Additional analyses are available in Appendix~\ref{sec:appExpLanguage}.

\paragraph{Scaling Laws.} Next, we examine power-law scaling properties, which are useful for predicting how performance will evolve as model size increases~\citep{Kaplan:20,Brown:20short}.
Figure~\ref{fig:temp_sclaw300B} illustrates the scaling trends across different architectures. Although each model exhibits comparable scaling patterns, their starting points vary. Among the tested architectures, RWKV-4 yields the lowest performance, with Llama and Mamba following, in ascending order. xLSTM outperforms Mamba by a similar margin as Mamba surpasses Llama. These results suggest that as model sizes grow, xLSTM is likely to maintain advantageous performance over both Transformer-based and State-Space approaches.

\textbf{Generation Times and Maximal Throughput.} In Figure~\ref{fig:inference_time_generation}, we report the duration required to generate text for various xLSTM variants at the 1.3B parameter scale, while Figure~\ref{fig:inference_time_decoding_throughput} (left) presents the highest attainable throughput. Our analysis includes comparable-sized implementations of Mamba, Llama, and RWKV taken from HuggingFace, where for the Llama model a static key-value cache is utilized. At the time when the experiments were conducted, both the complete compilation of cache for the Transformer Model and compiling Mamba with \texttt{torch.compile} were not operational. During text generation benchmarking, every model is assessed using batch size 1 and a pre-fill value of 16 tokens; this choice of pre-fill particularly favors the Transformer architecture.

From Figure~\ref{fig:inference_time_generation}, it becomes evident that both xLSTM and the other recurrent models, namely Mamba and RWKV-4, display linear growth with respect to input, in contrast to the quadratic growth characterizing Llama. To evaluate decoding throughput, we alter batch size and pre-fill settings in the Llama setup. Furthermore, as illustrated in Figure~\ref{fig:inference_time_decoding_throughput} (right), xLSTM's constant memory footprint allows it to support far larger batch sizes relative to Llama, resulting in the superior throughput achieved by xLSTM.

\section{Limitations}

(i) Unlike mLSTM, the sLSTM architecture restricts memory mixing in such a way that parallel processing cannot be achieved, thus limiting the possibility for efficient parallel computation. Despite this challenge, we created a rapid CUDA kernel for sLSTM, which currently runs at less than twice the duration required by our parallelized mLSTM implementation. (ii) The present CUDA kernels supporting mLSTM have not undergone optimization, which results in performance being approximately fourfold slower than either FlashAttention or the scan mechanism utilized in Mamba. Enhanced CUDA kernels, similar in approach to FlashAttention, may provide substantial speedup. (iii) The matrix memory component within mLSTM introduces significant computational demands because operations on $d \times d$ matrices are needed. Nonetheless, since memory access and updates are parameter-free and amenable to parallel execution via standard matrix processing techniques, the additional wall clock time resulting from the elaborate memory design remains marginal. (iv) Careful selection is required for the initialization parameters of the forget gates. (v) As the matrix memory does not scale with sequence length, increasing sequence length could potentially strain the memory capacity at longer context sizes. Empirically, this has not presented a limiting factor for contexts up to 16k tokens, as detailed in Section~\ref{sec:spaj300B_ctx_extrapolation}. (vi) Owing to the considerable computational resources required for training large language models, neither architectural nor hyperparameter optimization was performed extensively, particularly for bigger xLSTM variants. We expect that a thorough optimization campaign will be essential for xLSTM to achieve optimal performance.

\section{Conclusion}
We have provided a partial response to our fundamental inquiry: What is the potential of language modeling when scaling LSTM architectures to billions of parameters? To date, our findings suggest that such scaling achieves performance on par with prevailing approaches including Transformers and State Space Models. By introducing exponential gating combined with memory mixing and an innovative memory structure, we developed the xLSTM framework. When evaluated on language modeling tasks, xLSTM demonstrates superior or comparable results relative to leading models such as Transformers and State Space Models. Furthermore, the scaling behavior observed for xLSTM indicates that increasingly larger models could become substantial challengers to contemporary Large Language Models based on Transformer designs. The capabilities of xLSTM extend beyond language modeling and could offer notable advancements in domains such as Reinforcement Learning, Time Series Prediction, and modeling complex physical systems.

\end{document}